\documentclass[10pt, aspectratio=169]{beamer}
% Preámbulo modular para presentaciones Beamer (Modelación Estadística)
% Diseño y paleta institucional del Tecnológico de Monterrey

\documentclass[aspectratio=169,xcolor={dvipsnames,table}]{beamer}

\usepackage[utf8]{inputenc}
\usepackage[spanish,mexico]{babel}
\usepackage{amsmath,amssymb,amsfonts,mathtools}
\usepackage{graphicx}
\usepackage{listings}
\usepackage{booktabs}
\usepackage{tikz}
\usetikzlibrary{matrix,arrows,decorations.pathmorphing,shapes.geometric,positioning}

% Paleta Institucional Tecnológico de Monterrey
\definecolor{TecAzul}{HTML}{0039A6}       % Azul TEC: color primario institucional
\definecolor{TecAzulOscuro}{HTML}{1A2E51} % Azul marino oscuro
\definecolor{TecRojo}{HTML}{EC2661}       % Rojo de acento (paleta miTec)
\definecolor{TecGris}{HTML}{646464}       % Gris neutro para comentarios y subtítulos
\definecolor{TecFondoBloque}{HTML}{F4F6F9}% Fondo suave para bloques

% Configuración de Tema Beamer
\usetheme{Boadilla}
\usecolortheme[named=TecAzulOscuro]{structure}
\setbeamercolor{alerted text}{fg=TecRojo}
\setbeamercolor{palette primary}{bg=TecAzulOscuro,fg=white}
\setbeamercolor{palette secondary}{bg=TecAzul,fg=white}
\setbeamercolor{palette tertiary}{bg=TecAzulOscuro!80!black,fg=white}
\setbeamercolor{title}{fg=TecAzulOscuro,bg=TecFondoBloque}
\setbeamercolor{frametitle}{fg=TecAzulOscuro,bg=TecFondoBloque}

% Bloques personalizados elegantes
\setbeamercolor{block title}{bg=TecAzulOscuro,fg=white}
\setbeamercolor{block body}{bg=TecFondoBloque,fg=black}
\setbeamercolor{block title alerted}{bg=TecRojo,fg=white}
\setbeamercolor{block body alerted}{bg=TecRojo!8,fg=black}
\setbeamercolor{block title example}{bg=TecAzul,fg=white}
\setbeamercolor{block body example}{bg=TecAzul!8,fg=black}

% Redondear bordes y añadir sombra sutil
\setbeamertemplate{blocks}[rounded][shadow=true]
\setbeamertemplate{navigation symbols}{}

% Pie de página limpio con número de diapositiva
\setbeamertemplate{footline}{
  \leavevmode%
  \hbox{%
  \begin{beamercolorbox}[wd=.7\paperwidth,ht=2.25ex,dp=1ex,left]{author in head/foot}%
    \hspace*{2ex}\insertshortauthor~~(\insertshortinstitute) \hfill \insertshorttitle
  \end{beamercolorbox}%
  \begin{beamercolorbox}[wd=.3\paperwidth,ht=2.25ex,dp=1ex,right]{date in head/foot}%
    \insertshortdate{}\hspace*{2em}
    \insertframenumber{} / \inserttotalframenumber\hspace*{2ex}
  \end{beamercolorbox}}%
  \vskip0pt%
}

% Configuración de Listings de Python para visualización pedagógica de código
\lstset{
    language=Python,
    basicstyle=\ttfamily\footnotesize,
    keywordstyle=\color{TecAzulOscuro}\bfseries,
    stringstyle=\color{TecRojo},
    commentstyle=\color{TecGris}\itshape,
    showstringspaces=false,
    breaklines=true,
    frame=lines,
    rulecolor=\color{TecAzulOscuro!30},
    backgroundcolor=\color{TecFondoBloque},
    aboveskip=0.5em,
    belowskip=0.5em,
    literate={á}{{\'a}}1 {é}{{\'e}}1 {í}{{\'i}}1 {ó}{{\'o}}1 {ú}{{\'u}}1
             {Á}{{\'A}}1 {É}{{\'E}}1 {Í}{{\'I}}1 {Ó}{{\'O}}1 {Ú}{{\'U}}1
             {ñ}{{\~n}}1 {Ñ}{{\~N}}1 {¿}{{?`}}1 {¡}{{!`}}1
}

% Metadatos institucionales por defecto
\title[Modelación Estadística]{Modelación Estadística para Ciencia de Datos e Ingeniería}
\author[J. Castillo Colmenares]{Juliho David Castillo Colmenares}
\institute[Tec de Monterrey]{Tecnológico de Monterrey \\ Escuela de Ingeniería y Ciencias}
\date{\today}

% Comandos y macros estadísticas compartidas para las presentaciones Beamer (Metropolis)

% Conjuntos numéricos básicos
\newcommand{\R}{\mathbb{R}}
\newcommand{\N}{\mathbb{N}}
\newcommand{\Q}{\mathbb{Q}}
\newcommand{\Z}{\mathbb{Z}}
\newcommand{\set}[1]{\left\{ #1 \right\}}
\newcommand{\abs}[1]{\left|#1\right|}
\newcommand{\norm}[1]{\left\| #1 \right\|}

% Comandos estadísticos y probabilísticos del curso
\newcommand{\Var}{\operatorname{Var}}
\newcommand{\cov}{\operatorname{Cov}}
\newcommand{\corr}{\rho}
\newcommand{\comb}[2]{\begin{pmatrix} #1 \\ #2 \end{pmatrix}}
\newcommand{\s}{\sigma}
\newcommand{\particion}{\mathfrak{P}}
\newcommand{\card}[1]{\#\left( #1 \right)}
\newcommand{\seq}[3]{\set{#1}_{#2}^{#3}}
\newcommand{\E}{\mathbb{E}}
\newcommand{\Prob}{\mathbb{P}}

% Entornos pedagógicos y cajas en español académico natural
\newenvironment{discusion}{%
  \begin{alertblock}{Pregunta de Discusión}%
}{%
  \end{alertblock}%
}

\newenvironment{reflexion}{%
  \begin{alertblock}{Reflexión para el Aula}%
}{%
  \end{alertblock}%
}

\newenvironment{paradoja}{%
  \begin{alertblock}{Paradoja de Exploración}%
}{%
  \end{alertblock}%
}

\newenvironment{motivacion}{%
  \begin{exampleblock}{Motivación: Ciencia de Datos en la Práctica}%
}{%
  \end{exampleblock}%
}

\newenvironment{retoclase}{%
  \begin{block}{Reto Rápido en Equipo}%
}{%
  \end{block}%
}


\title[Regresión con scikit-learn]{Sección 09.10: Regresión Lineal con \texttt{scikit-learn}}
\subtitle{El Flujo de Trabajo Profesional y la Selección Automática de Variables}
\author[J. Castillo Colmenares]{Juliho Castillo Colmenares}
\institute{Tecnológico de Monterrey}
\date{\vspace{-1.2cm}}

\begin{document}

% Slide 1: Portada
\begin{frame}[plain]
  \titlepage
\end{frame}

% Slide 2: Hoja de Ruta
\begin{frame}{Hoja de Ruta --- Capítulo 09: Regresiones Lineales y Múltiples}
  \scriptsize
  ¿Dónde nos encontramos en el desarrollo modular del capítulo?
  \begin{itemize}\itemsep=0.02em
    \item \textbf{09.01} Correlación como Premisa de la Regresión
    \item \textbf{09.02} Introducción a la Regresión Lineal
    \item \textbf{09.03} Matemáticas de la Regresión: Derivación MCO y $R^2$
    \item \textbf{09.04} Regresión Lineal sobre Datos Simulados
    \item \textbf{09.05} Coeficientes Óptimos, Pruebas $t$/$F$ y RSE
    \item \textbf{09.06} Implementación en Python con \texttt{statsmodels}
    \item \textbf{09.07} Regresión Múltiple, Ecuación Normal y Ridge/Lasso
    \item \textbf{09.08} Validación de Modelos y $k$-fold Cross-Validation
    \item \textbf{09.09} Diagnóstico de Residuos y Multicolinealidad (VIF)
    \item \textbf{09.10} Regresión con \texttt{scikit-learn} \emph{(Hoy)}
    \item \textbf{09.11} Variables Categóricas y Variables Muda
    \item \textbf{09.12} Transformaciones No Lineales y Regresión Polinomial
  \end{itemize}
  \pause
  \vspace{-0.05cm}
  \begin{block}{Objetivo de la Sesión}
    Implementar regresión lineal con \texttt{scikit-learn}, el estándar de la industria en aprendizaje automático; verificar que resuelve exactamente la misma Ecuación Normal derivada en la Sección 09.07 mediante descomposición en valores singulares (SVD); y automatizar la selección de variables con Eliminación Recursiva de Características (RFE).
  \end{block}
\end{frame}

% Slide 3: Motivación
\begin{frame}{Nota Metodológica: la Única Excepción del Capítulo}
  \footnotesize
  \begin{motivacion}
    Todo el capítulo ha construido sus laboratorios exclusivamente con \texttt{numpy}/\texttt{scipy}, incluso para reproducir lo que \texttt{statsmodels} calcula (Sección 09.06). Esta sección es la \textbf{única excepción}: su tema \emph{es} la librería \texttt{scikit-learn} en sí misma, así que la usamos directamente.
  \end{motivacion}

  \vspace{0.15cm}
  \pause
  \texttt{scikit-learn} adopta un enfoque de \emph{aprendizaje automático} (interfaz \texttt{fit}/\texttt{predict}/\texttt{score} uniforme para todos los modelos), distinto al enfoque \emph{inferencial} de \texttt{statsmodels} (tablas de resumen estadístico).
\end{frame}

% Slide 4: El flujo fit/predict/score
\begin{frame}[fragile]{El Flujo de Trabajo \texttt{fit}/\texttt{predict}/\texttt{score}}
  \footnotesize
  \begin{block}{Ajuste}
    \begin{verbatim}
X_train, X_test, y_train, y_test = train_test_split(X, y, test_size=0.2)
lm = LinearRegression()
lm.fit(X_train, y_train)
    \end{verbatim}
  \end{block}
  \pause

  \vspace{0.1cm}
  \begin{alertblock}{Predicción y Evaluación}
    \texttt{lm.coef\_} y \texttt{lm.intercept\_} contienen exactamente $\hat{\boldsymbol\beta}$; \texttt{lm.predict(X\_test)} calcula $\hat{\mathbf{Y}}$; \texttt{lm.score(X,y)} retorna $R^2$ sobre el conjunto proporcionado.
  \end{alertblock}
\end{frame}

% Slide 5: SVD vs Ecuación Normal
\begin{frame}{¿Cómo Resuelve \texttt{scikit-learn} el Sistema?}
  \footnotesize
  \begin{block}{No Invierte $\mathbf{X}^T\mathbf{X}$ Directamente}
    Internamente, \texttt{LinearRegression} usa la descomposición en valores singulares (SVD) de $\mathbf{X}=\mathbf{U}\boldsymbol\Sigma\mathbf{V}^T$, resolviendo $\hat{\boldsymbol\beta}=\mathbf{X}^+\mathbf{Y}$ vía la pseudoinversa de Moore-Penrose $\mathbf{X}^+=\mathbf{V}\boldsymbol\Sigma^{-1}\mathbf{U}^T$.
  \end{block}
  \pause

  \vspace{0.1cm}
  \begin{alertblock}{Equivalencia y Ventaja Numérica}
    Cuando $\mathbf{X}$ tiene rango completo, esto es algebraicamente idéntico a $(\mathbf{X}^T\mathbf{X})^{-1}\mathbf{X}^T\mathbf{Y}$ (Sección 09.07). La SVD evita elevar al cuadrado el número de condición de $\mathbf{X}$, siendo mucho más estable ante multicolinealidad severa (Sección 09.09).
  \end{alertblock}
\end{frame}

% Slide 6: RFE
\begin{frame}[fragile]{Selección Automática: Eliminación Recursiva de Características}
  \footnotesize
  \begin{block}{\texttt{RFE} (\emph{Recursive Feature Elimination})}
    \begin{verbatim}
selector = RFE(LinearRegression(), n_features_to_select=2)
selector.fit(X, Y)
    \end{verbatim}
    Ajusta el modelo con todas las variables, elimina la de menor peso, y repite hasta dejar el número deseado.
  \end{block}
  \pause

  \vspace{0.1cm}
  \begin{alertblock}{Salida}
    \texttt{selector.support\_} indica qué variables sobreviven; \texttt{selector.ranking\_} ordena las eliminadas por importancia relativa (rango 1 = seleccionada).
  \end{alertblock}
\end{frame}

% Slide 7: Lab Python (Bloque 1)
\begin{frame}[fragile]{Laboratorio en Python: Flujo de Trabajo de \texttt{scikit-learn}}
  \scriptsize
  Ajuste con \texttt{train\_test\_split} + \texttt{LinearRegression} (\texttt{09.10\_scikit\_learn\_regression.py}):
  {\lstset{basicstyle=\fontsize{3.2pt}{3.9pt}\selectfont\ttfamily,aboveskip=0.05em,belowskip=0.05em}
  \lstinputlisting[firstline=24, lastline=44]{../../code/09_regresiones/09.10_scikit_learn_regression.py}}
\end{frame}

% Slide 8: Lab Python (Bloque 2)
\begin{frame}[fragile]{Laboratorio en Python: Verificación contra MCO Manual}
  \scriptsize
  Comparación exacta entre \texttt{scikit-learn} y la Ecuación Normal de la Sección 09.07:
  {\lstset{basicstyle=\fontsize{3.4pt}{4.1pt}\selectfont\ttfamily,aboveskip=0.05em,belowskip=0.05em}
  \lstinputlisting[firstline=47, lastline=59]{../../code/09_regresiones/09.10_scikit_learn_regression.py}}
\end{frame}

% Slide 9: Lab Python (Bloque 3)
\begin{frame}[fragile]{Laboratorio en Python: Eliminación Recursiva de Características}
  \scriptsize
  Selección automática de las 2 variables más relevantes entre 3 candidatas:
  {\lstset{basicstyle=\fontsize{3.0pt}{3.7pt}\selectfont\ttfamily,aboveskip=0.05em,belowskip=0.05em}
  \lstinputlisting[firstline=62, lastline=82]{../../code/09_regresiones/09.10_scikit_learn_regression.py}}
\end{frame}

% Slide 10: Salida Terminal
\begin{frame}[fragile]{Laboratorio en Python: Salida en Terminal}
  \scriptsize
  Salida estándar generada al ejecutar el script del laboratorio en Python:
  \vspace{0.02cm}
  \begin{block}{Salida Estándar en Consola}
    \fontsize{3.4pt}{4.1pt}\selectfont
    \begin{verbatim}
=== Block 1: Fit, Score, and Predict ===
Intercept: 2.8772, Coefficients (TV, Radio): [0.0455 0.1947]
R^2 train=0.8633, R^2 test=0.9016

=== Block 2: Cross-Checking Against Scratch OLS ===
From-scratch: [2.877152 0.04554  0.194653]
scikit-learn: [2.877152 0.04554  0.194653]
Max abs difference: 1.95e-14

=== Block 3: RFE ===
Support: TV=True, Radio=True, Newspaper=False
Ranking: TV=1, Radio=1, Newspaper=2
    \end{verbatim}
  \end{block}
\end{frame}

% Slide 11A: Ejercicio Nivel 1 (Enunciado)

% Slide 11B: Ejercicio Nivel 1 (Resolución)

% Slide 12A: Ejercicio Nivel 2 (Enunciado)

% Slide 12B: Ejercicio Nivel 2 (Resolución)

% Slide 13A: Ejercicio Nivel 3 (Enunciado)
\begin{frame}{Ejercicio en Clase --- Nivel Analítico (1/2)}
  \footnotesize
  \begin{block}{Problema --- SVD vs. Ecuación Normal (Enunciado, Problema 9.17.7)}
    Explique por qué la SVD y la inversión directa de $\mathbf{X}^T\mathbf{X}$ dan el mismo resultado cuando $\mathbf{X}$ tiene rango completo, y por qué la SVD es numéricamente superior ante multicolinealidad severa.
  \end{block}

  \vspace{0.1cm}
  \pause
  \begin{alertblock}{Estrategia}
    Compare el número de condición de $\mathbf{X}$ contra el de $\mathbf{X}^T\mathbf{X}$.
  \end{alertblock}
\end{frame}

% Slide 13B: Ejercicio Nivel 3 (Resolución)
\begin{frame}{Ejercicio en Clase --- Nivel Analítico (2/2)}
  \scriptsize
  El número de condición de $\mathbf{X}^T\mathbf{X}$ es el cuadrado del de $\mathbf{X}$: \pause

  \vspace{0.05cm}
  \begin{block}{Cálculo}
    \scriptsize
    $$\kappa(\mathbf{X}^T\mathbf{X}) = \kappa(\mathbf{X})^2$$
  \end{block}

  \vspace{0.05cm}
  \pause
  \begin{alertblock}{Interpretación}
    \scriptsize
    Invertir $\mathbf{X}^T\mathbf{X}$ directamente amplifica cuadráticamente los errores de redondeo cuando $\mathbf{X}$ ya está mal condicionada por multicolinealidad. La SVD opera sobre $\mathbf{X}$ directamente, evitando ese cuadrado y produciendo una solución mucho más estable.
  \end{alertblock}
\end{frame}

% Slide 14A: Ejercicio Nivel 4 (Enunciado)
\begin{frame}{Ejercicio en Clase --- Nivel Desafiante (1/2)}
  \footnotesize
  \begin{block}{Problema --- Equivalencia SVD--Ecuación Normal (Enunciado, Problema 9.17.9)}
    Con $\mathbf{X}=\mathbf{U}\boldsymbol\Sigma\mathbf{V}^T$, demuestre que $\hat{\boldsymbol\beta}=\mathbf{X}^+\mathbf{Y}=\mathbf{V}\boldsymbol\Sigma^{-1}\mathbf{U}^T\mathbf{Y}$ es idéntico a $(\mathbf{X}^T\mathbf{X})^{-1}\mathbf{X}^T\mathbf{Y}$ cuando $\mathbf{X}$ tiene rango completo.
  \end{block}

  \vspace{0.1cm}
  \pause
  \begin{alertblock}{Estrategia}
    Sustituya la SVD en $\mathbf{X}^T\mathbf{X}$ y en $\mathbf{X}^T\mathbf{Y}$ usando la ortogonalidad de $\mathbf{U},\mathbf{V}$.
  \end{alertblock}
\end{frame}

% Slide 14B: Ejercicio Nivel 4 (Resolución)
\begin{frame}{Ejercicio en Clase --- Nivel Desafiante (2/2)}
  \scriptsize
  Sustituimos y simplificamos usando $\mathbf{U}^T\mathbf{U}=\mathbf{V}^T\mathbf{V}=\mathbf{I}$: \pause

  \vspace{0.05cm}
  \begin{block}{Cálculo}
    \scriptsize
    $$\mathbf{X}^T\mathbf{X}=\mathbf{V}\boldsymbol\Sigma^2\mathbf{V}^T \implies (\mathbf{X}^T\mathbf{X})^{-1}=\mathbf{V}\boldsymbol\Sigma^{-2}\mathbf{V}^T$$
    $$\hat{\boldsymbol\beta}=\mathbf{V}\boldsymbol\Sigma^{-2}\mathbf{V}^T(\mathbf{V}\boldsymbol\Sigma\mathbf{U}^T\mathbf{Y})=\mathbf{V}\boldsymbol\Sigma^{-1}\mathbf{U}^T\mathbf{Y}=\mathbf{X}^+\mathbf{Y}$$
  \end{block}

  \vspace{0.05cm}
  \pause
  \begin{alertblock}{Interpretación}
    Ambas expresiones son algebraicamente idénticas: la SVD no introduce matemática nueva, solo una ruta numéricamente más estable hacia la misma solución.
  \end{alertblock}
\end{frame}

% Slide 15: Cierre
\begin{frame}{Cierre de Sección: Regresión con \texttt{scikit-learn}}
  \begin{reflexion}
    \texttt{scikit-learn} no introduce matemática nueva: resuelve exactamente la misma Ecuación Normal derivada en la Sección 09.07, solo que mediante SVD por robustez numérica, y automatiza la selección de variables que antes hacíamos manualmente comparando modelos (Secciones 09.02, 09.07-09.08).
  \end{reflexion}

  \vspace{0.2cm}
  \pause
  \begin{block}{Lecciones Clave}
    \begin{itemize}\itemsep=0.08cm
      \item \textbf{\texttt{fit}/\texttt{predict}/\texttt{score}:} interfaz uniforme del ecosistema de aprendizaje automático.
      \item \textbf{SVD $\equiv$ Ecuación Normal:} misma matemática, mayor estabilidad numérica.
      \item \textbf{RFE:} automatiza la selección de variables mediante eliminación recursiva basada en importancia.
    \end{itemize}
  \end{block}
\end{frame}

% Slide 16: Perspectiva Modular
\begin{frame}{Perspectiva Modular --- El Próximo Reto}
  \scriptsize
  \begin{retoclase}
    Hasta ahora todos los predictores han sido numéricos. Pero muchos conjuntos de datos reales incluyen variables \textbf{categóricas} (género, categoría de producto, región). La Sección 09.11 muestra cómo codificarlas correctamente mediante variables muda (\emph{dummy variables}) para incorporarlas en un modelo de regresión.
  \end{retoclase}

  \vspace{0.08cm}
  \pause
  \begin{alertblock}{Preparación Recomendada}
    \begin{itemize}\itemsep=0.06cm
      \item Reflexiona sobre por qué asignar códigos numéricos arbitrarios (0, 1, 2...) a categorías sin orden natural sería matemáticamente incorrecto.
      \item Repasa la función \texttt{pandas.get\_dummies} si no la conoces.
    \end{itemize}
  \end{alertblock}
\end{frame}

\end{document}
